\documentclass{subfiles}
\begin{document}
    Projective spaces are deeply rooted in art,
        as these were implicitly used when dealing with perspective.
        As it turns out,
        perspective is nothing else but a consequence of how the projective space works.

    Formally speaking, an \(n\)-dimensional projective space over a field \(\K\) is defined 
        as the set of non-empty \emph{rays} \(\mathbf{v} \in \K^{n + 1}\).
        More precisely 
        \[
            \mathbb{P}^{n}(\K) = \set{P = \iprod{\mathbf{v}}}[\mathbf{0} \ne \mathbf{v}  \in \K^{n + 1}],
        \]
        where \(\mathbf{v} = [x_{0}, \ldots, x_{n}], x_{i} \in \K\).

    Points in \(\mathbb{P}^{n}(\K)\) are represented using two types of coordinates:
        \begin{enumerate}
            \item \emph{Homogeneous} coordinates, by which a point \(Q\) differs from \(P\)
                if and only if \(Q \ne \lambda P\). In other words,
                using homogeneous coordinates,
                all points that differ by a multiplicative constant represent the same point.

            \item \emph{Non-homogeneous} coordinates, by which \(P\) is written as 
                \[ P = \begin{cases}
                    [1, \sfrac{x_{1}}{x_{0}}, \ldots, \sfrac{x_{n}}{x_{0}}], \text{ if } x_{0} \ne 0; \\
                    [0, x_{1}, \ldots, x_{n}], \text{ if } x_{0} = 0.
                \end{cases}\]
        \end{enumerate}
        When using non-homogeneous coordinates, 
        we say that \(P = [x_{0}, \ldots, x_{n}]\) is a proper point, or equivalently,
        an affine point if \(x_{0} \ne 0\); we say that \(P\) is an improper point,
        or equivalently, a point at infinity if \(x_{0} = 0\).

    Our discussion on elliptic curves will involve points in \(\mathbb{P}^{2}(\K)\),
        often represented by non-homogeneous coordinates.
        For this reason, to ease the notation a bit, 
        we write \(P = [1, \sfrac{x_{1}}{x_{0}}, \sfrac{x_{2}}{x_{0}}]\) as
        \(P = [1, x, y]\), where \(x = \sfrac{x_{1}}{x_{0}}\) and \(y = \sfrac{x_{2}}{x_{0}}\).
        For point at infinity we write \(P\) as \(P = [0, X, Y]\).

    \subsubsection{Homogeneous polynomials}
    \subfile{../subsub/3.1.1 - Homogeneous polynomials}

    \subsubsection{Elliptic curves}
    \subfile{../subsub/3.1.2 - Elliptic curves}
        
\end{document}
